\documentclass[light,minted=false]{cuzbeamer} % tema claro, sin resaltado de código

\usepackage{graphicx}

\title{Aplicación de Inteligencia Artificial Generativa en la Visualización de Datos}
\subtitle{Con Python, Plotly y Dash}
\author{Dr. Alejandro González Turrubiates}
\institute{Universidad Autónoma de Tamaulipas}
\date{2 al 4 de septiembre de 2026}
\email{agturrubiates@uat.edu.mx}

% Referencia textual del congreso donde se imparte el taller (el logo del
% evento ya ocupa, en su lugar, la esquina superior derecha de cada slide:
% ver images/logo_SIP_UAT-trimmed.png y images/logo_SIP_UAT-negro.png).
\AtBeginDocument{%
    \addtobeamertemplate{title page}{}{%
        \begin{tikzpicture}[remember picture,overlay]
            \node[anchor=south west,xshift=0.3cm,yshift=0.35cm,align=left,font=\tiny,text=white]
                at (current page.south west)
                {En el marco del Congreso Internacional CIREDII 2026 --- Colima};
        \end{tikzpicture}%
    }%
}

\begin{document}
    \maketitle

    \begin{frame}{Objetivos del taller}
        \begin{itemize}
            \item Formular prompts efectivos para generar visualizaciones
            \item Generar y ajustar gráficas a partir de datos reales
            \item \textbf{Verificar} los resultados generados por IA
            \item Aplicar buenas prácticas de uso responsable de la IA
        \end{itemize}
    \end{frame}

    \begin{frame}{Agenda del día (9:00--14:00)}
        \begin{enumerate}
            \item Fundamentos de prompting
            \item Plotly guiado
            \item Práctica libre
            \item De Plotly a Dash
            \item Buenas prácticas y cierre
        \end{enumerate}
    \end{frame}

    \begin{frame}{Antes de empezar}
        Verifica que tienes instalado:
        \begin{itemize}
            \item Python 3.11+
            \item Una CLI agentic (Gemini CLI o Claude Code)
        \end{itemize}

        \vspace{1em}
        Ver \texttt{requisitos/} si algo falta.
    \end{frame}
\end{document}
